\chapter*{Chapitre 2 — Hydraulique des circuits}
\addcontentsline{toc}{chapter}{Chapitre 2 — Hydraulique des circuits}

% ============================================================
\section*{Énoncés}
% ============================================================

\exercice{Pertes de charge d'un réseau bitubes}

Un réseau bitubes de chauffage alimente 3 radiateurs en dérivation. Le tronçon commun (départ chaudière → nœud de distribution) est en tube acier DN 32 (diamètre intérieur \SI{35{,}9}{\milli\metre}), longueur \SI{20}{\metre}. Il débite \SI{1500}{\litre\per\hour}. La perte de charge linéique pour ce tronçon est $R = \SI{180}{\pascal\per\metre}$. On majore de 40\,\% pour les singularités.

\begin{enumerate}
  \item Calculer la perte de charge du tronçon commun.
  \item Calculer la vitesse du fluide. Vérifier qu'elle est dans la plage recommandée.
  \item La branche la plus défavorisée (radiateur le plus éloigné) présente une perte de charge de \SI{9000}{\pascal}. Les deux autres branches ont respectivement \SI{4500}{\pascal} et \SI{6000}{\pascal}. Calculer la perte de charge supplémentaire à créer sur chaque robinet d'équilibrage.
\end{enumerate}

% ----------------------------------------------------------

\exercice{Dimensionnement d'une installation de climatisation eau glacée}

Un bâtiment tertiaire nécessite une puissance de froid de \SI{120}{\kilo\watt}. Le groupe froid produit de l'eau glacée à 7\,°C retournant à 12\,°C ($\Delta T = 5$\,°C, $\rho = \SI{999}{\kilogram\per\cubic\metre}$, $c_p = \SI{4182}{\joule\per\kilogram\per\kelvin}$). La perte de charge de la boucle la plus défavorisée est estimée à \SI{35000}{\pascal}.

\begin{enumerate}
  \item Calculer le débit volumique nécessaire en m³/h.
  \item Calculer la HMT de la pompe en mCE.
  \item Calculer la puissance hydraulique absorbée (en supposant $\eta_{pompe} = 55\,\%$).
\end{enumerate}

% ----------------------------------------------------------

\exercice{Analyse d'un dysfonctionnement hydraulique}

Lors de la réception d'une installation de chauffage, un technicien mesure les débits suivants sur 3 radiateurs d'un même niveau :

\begin{center}
\begin{tabular}{lccc}
  \toprule
  Radiateur & Débit mesuré (L/h) & Débit nominal (L/h) & Écart \\
  \midrule
  R1 (proche pompe) & 280 & 180 & +56\,\% \\
  R2 (intermédiaire) & 195 & 180 & +8\,\% \\
  R3 (éloigné) & 120 & 180 & -33\,\% \\
  \bottomrule
\end{tabular}
\end{center}

\begin{enumerate}
  \item Identifier le problème et sa cause probable.
  \item Proposer les actions correctives à mettre en œuvre.
  \item Calculer la perte de charge supplémentaire à créer sur R1 sachant que R3 est la branche index avec une $\Delta P$ mesurée de \SI{7200}{\pascal}, et que la $\Delta P$ sur R1 est de \SI{4100}{\pascal}.
\end{enumerate}

\newpage
% ============================================================
\section*{Corrigés}
% ============================================================

\subsection*{Corrigé — Exercice 1}

\textit{Question 1 — Perte de charge tronçon commun :}
\[
  \Delta P = R \times L \times (1 + 0{,}40) = 180 \times 20 \times 1{,}40 = \SI{5040}{\pascal}
\]

\textit{Question 2 — Vitesse du fluide :}
\[
  S = \frac{\pi \times (0{,}0359)^2}{4} = \SI{1{,}012e-3}{\square\metre}
\]
\[
  q_v = \frac{1500}{3600 \times 1000} = \SI{4{,}17e-4}{\cubic\metre\per\second}
\]
\[
  v = \frac{4{,}17 \times 10^{-4}}{1{,}012 \times 10^{-3}} = \SI{0{,}41}{\metre\per\second}
\]
Vitesse dans la plage recommandée (0,5 à 1,5 m/s) — légèrement en dessous, acceptable pour une colonne montante.

\textit{Question 3 — Équilibrage :}

Branche index = \SI{9000}{\pascal}
\begin{itemize}
  \item Branche 2 : excédent = $9000 - 6000 = \SI{3000}{\pascal}$ → régler le robinet à +3000\,Pa
  \item Branche 3 : excédent = $9000 - 4500 = \SI{4500}{\pascal}$ → régler le robinet à +4500\,Pa
\end{itemize}

% ----------------------------------------------------------

\subsection*{Corrigé — Exercice 2}

\textit{Question 1 — Débit :}
\[
  \dot{m} = \frac{120\,000}{4182 \times 5} = \SI{5{,}74}{\kilogram\per\second}
\]
\[
  q_v = \frac{5{,}74}{999} = \SI{5{,}74e-3}{\cubic\metre\per\second} = \SI{20{,}67}{\cubic\metre\per\hour}
\]

\[
  \boxed{q_v \approx \SI{20{,}67}{\cubic\metre\per\hour}}
\]

\textit{Question 2 — HMT :}
\[
  HMT = \frac{35\,000}{999 \times 9{,}81} = \SI{3{,}57}{mCE}
\]

\[
  \boxed{HMT \approx \SI{3{,}57}{mCE}}
\]

\textit{Question 3 — Puissance absorbée :}
\[
  P_{hyd} = \rho \cdot g \cdot q_v \cdot HMT = 999 \times 9{,}81 \times 5{,}74 \times 10^{-3} \times 3{,}57 = \SI{201}{\watt}
\]
\[
  P_{abs} = \frac{P_{hyd}}{\eta} = \frac{201}{0{,}55} = \SI{365}{\watt} \approx \SI{0{,}37}{\kilo\watt}
\]

\[
  \boxed{P_{abs} \approx \SI{0{,}37}{\kilo\watt}}
\]

% ----------------------------------------------------------

\subsection*{Corrigé — Exercice 3}

\textit{Question 1 :} Le réseau est clairement déséquilibré : R1 reçoit trop de débit (surchauffe), R3 est sous-alimenté (inconfort froid). Cause : absence d'équilibrage hydraulique ou robinets d'équilibrage non réglés.

\textit{Question 2 :} Actions correctives :
\begin{itemize}
  \item Identifier et vérifier les robinets d'équilibrage sur chaque branche
  \item Utiliser une pince de mesure de débit ou un débitmètre ultrasonique
  \item Régler les robinets d'équilibrage selon la méthode de la branche index (R3)
  \item Re-mesurer les débits après réglage et valider
\end{itemize}

\textit{Question 3 :}
\[
  \Delta P_{excédent} = \Delta P_{R3} - \Delta P_{R1} = 7200 - 4100 = \SI{3100}{\pascal}
\]

\[
  \boxed{\Delta P_{excédent} = \SI{3100}{\pascal}}
\]

Le robinet d'équilibrage de R1 doit être réglé pour créer une perte de charge supplémentaire de \SI{3100}{\pascal}.
